\chapter{État de l'art}
\chapminitoc
% Chapitre « niveau PFE » : montrer que tu connais le domaine AVANT de montrer
% ce que tu as fait. Chaque section : 0,5 à 1,5 page + références en biblio.

\section{Le véhicule autonome et connecté}
% TODO : niveaux d'automatisation SAE (0 à 5), architecture
% perception -> décision -> action, enjeux sécurité/société.

\section{Les middlewares robotiques : ROS 2 et DDS}
% TODO : pourquoi ROS 2 (temps réel, QoS, industrialisation), principe DDS,
% comparaison FastDDS vs Zenoh (découverte, NAT, réseaux mobiles) — ce
% comparatif justifiera tes choix du chapitre 5.

\section{La téléopération de véhicules et l'apport de la 5G}
% TODO : exigences de latence pour la conduite à distance (~<100 ms),
% architectures 5G (NSA/SA), cas d'usage industriels existants.

\section{Les communications V2X}
% TODO : ITS-G5 (802.11p) vs C-V2X ; normes ETSI : CAM (EN 302 637-2),
% DENM (EN 302 637-3), GeoNetworking ; piles logicielles dont Vanetza.
% C'est le socle théorique de TON chapitre 7.

\section{Perception embarquée}
% TODO : lidar 2D, détection par apprentissage profond (famille YOLO),
% fusion capteurs, principe de l'AEB.

\section{Localisation et navigation autonome}
% TODO : SLAM (slam_toolbox), pile Nav2, méthodes réactives (follow the gap,
% issue de la compétition F1TENTH).

\section{Les plateformes à échelle réduite}
% TODO : F1TENTH, Duckietown, etc. : pourquoi la recherche utilise des
% voitures 1/10 — coût, sécurité, itération rapide. Positionner ProtoVA.

\section{Synthèse : les verrous adressés par ce stage}
% TODO : 4-5 verrous concrets qui font la transition vers le chapitre 3
% (ex. « la latence 5G est-elle compatible avec la téléopération ? »,
% « les piles ETSI sont-elles utilisables sans matériel 802.11p dédié ? »).
